% ─────────────────────────────────────────────────────────────────────────────
% Statistical Significance Subsection — correct tests per metric
% ─────────────────────────────────────────────────────────────────────────────

\subsubsection{Statistical Significance}
\label{sec:statistical_significance}

Not all metrics reported in this evaluation are amenable to formal hypothesis
testing, so we first clarify which metrics warrant a test and why.

\textbf{M1 (Reaction Performance)} is defined as the ratio of post-novelty to
pre-novelty cumulative action cost. Because the pre-novelty value is the
normalisation denominator, the default M1 is always exactly $1.00$ with zero
variance by construction. Comparing a distribution to a constant violates the
assumptions of any two-sample test; M1 is therefore interpreted descriptively
and no significance test is applied.

\textbf{M3 (Detection Rate)} is a binary, deterministic outcome. The planning
agent achieved a $100\%$ detection rate in every novelty scenario, leaving no
distributional uncertainty to test.

\textbf{M2 (Steps)} is the only planning-agent metric with genuine variance on
both sides of the novelty boundary. We treat the pre-novelty mean as a fixed
reference value and apply a \textbf{one-sample $t$-test} to the five
post-novelty observations, asking whether the post-novelty step count differs
significantly from the pre-novelty baseline.

\textbf{RL success rate} is computed from $n = 10$ independent trials, each a
genuinely separate experimental run with a different random seed. We apply an
\textbf{unpaired Welch's $t$-test}~\cite{kim2015t} between the default and each
post-novelty condition, with a significance level of $\alpha = 0.05$.

%
% ── M2: PLANNING AGENT ───────────────────────────────────────────────────────
%
\begin{table}[!tb]
\footnotesize
\begin{center}
\begin{tabular}{>{\centering}p{3.2cm}>{\centering}p{2.2cm}>{\centering}p{1.2cm}>{\centering}p{1.8cm}p{0.1cm}}

\arrayrulecolor{taupegray}\toprule
\textbf{\begin{small}Novelty\end{small}} &
\textbf{\begin{small}Pre-novelty (ref.)\end{small}} &
{\textbf{\begin{small}$p$-value\end{small}}} &
\textbf{\begin{small}Statistically Significant\end{small}}\tabularnewline
\midrule

\multicolumn{4}{c}{\textbf{NovelGridworld-Bow-v0}}\tabularnewline
\midrule
Fence        & $87.6 \pm 16.8$ & $0.0117$ & Yes \tabularnewline
Arrow/Spring & $87.6 \pm 16.8$ & $0.6798$ & No  \tabularnewline
Axe          & $87.6 \pm 16.8$ & $0.2174$ & No  \tabularnewline
\midrule

\multicolumn{4}{c}{\textbf{NovelGridworld-Bow-v1}}\tabularnewline
\midrule
Fence        & $36.8 \pm 6.5$ & $0.0013$ & Yes \tabularnewline
Arrow/Spring & $36.8 \pm 6.5$ & $0.9380$ & No  \tabularnewline
Axe          & $36.8 \pm 6.5$ & $0.4776$ & No  \tabularnewline
\midrule

\multicolumn{4}{c}{\textbf{NovelGridworld-Pogostick-v0}}\tabularnewline
\midrule
Fence        & $127.8 \pm 9.0$ & $0.5482$ & No  \tabularnewline
Arrow/Spring & $127.8 \pm 9.0$ & $0.0341$ & Yes \tabularnewline
Axe          & $127.8 \pm 9.0$ & $0.0019$ & Yes \tabularnewline
\midrule

\multicolumn{4}{c}{\textbf{NovelGridworld-Pogostick-v1}}\tabularnewline
\midrule
Fence        & $94.2 \pm 12.9$ & $0.1651$ & No  \tabularnewline
Arrow/Spring & $94.2 \pm 12.9$ & $0.5490$ & No  \tabularnewline
Axe          & $94.2 \pm 12.9$ & $0.2525$ & No  \tabularnewline

\bottomrule
\multicolumn{4}{l}{\small One-sample $t$-test, $n=5$ post-novelty episodes, $\alpha=0.05$.}\\
\multicolumn{4}{l}{\small Pre-novelty mean treated as fixed reference value.}\\
\end{tabular}
\caption{\small Statistical significance of the change in M2 (steps) for the
planning agent across all tasks and novelty scenarios.}
\label{tab:stat_sig_planning}
\end{center}
\end{table}

Table~\ref{tab:stat_sig_planning} shows that the Fence novelty produces a
significant increase in step count in the Bow environments, confirming that the
detrimental effect of fence objects is not attributable to sampling variability.
The result is not significant in the Pogostick environments, where higher
step-count variance reduces the power of the test.
The Arrow/Spring and Axe novelties are largely non-significant across all tasks,
consistent with the expectation that an irrelevant novelty leaves performance
unchanged and the beneficial effect of the axe is too small to detect reliably
at $n = 5$.
The significant results observed for Arrow/Spring and Axe in Pogostick-v0 are
noteworthy; since these novelties are expected to have little effect, these
results likely reflect the small sample size ($n = 5$) rather than a genuine
performance change, and should be interpreted with caution.

%
% ── RL AGENT ─────────────────────────────────────────────────────────────────
%
\begin{table}[!tb]
\footnotesize
\begin{center}
\begin{tabular}{>{\centering}p{4.2cm}>{\centering}p{1.2cm}>{\centering}p{1.8cm}p{0.1cm}}

\arrayrulecolor{taupegray}\toprule
\textbf{\begin{small}Comparison (Default $\leftrightarrow$ Novelty)\end{small}} &
{\textbf{\begin{small}$p$-value\end{small}}} &
\textbf{\begin{small}Statistically Significant\end{small}}\tabularnewline
\midrule

\multicolumn{3}{c}{\textbf{NovelGridworld-Bow-v0}}\tabularnewline
\midrule
Default $\leftrightarrow$ Action remapping & $0.4146$   & No  \tabularnewline
Default $\leftrightarrow$ Fire wall        & $< 0.0001$ & Yes \tabularnewline
Default $\leftrightarrow$ Break increase   & $< 0.0001$ & Yes \tabularnewline

\bottomrule
\multicolumn{3}{l}{\small Unpaired Welch's $t$-test, $n = 10$ independent trials, $\alpha = 0.05$.}\\
\end{tabular}
\caption{\small Statistical significance of the change in success rate for the
RL agent on \emph{NovelGridworld-Bow-v0}. The non-significant result for Action
remapping is partly attributable to high variance in that condition
($\sigma = 40.0\%$), which limits the power of the test.}
\label{tab:stat_sig_rl}
\end{center}
\end{table}

Table~\ref{tab:stat_sig_rl} confirms that the Fire-wall and Break-increase
novelties cause a statistically significant drop in task success rate. The
Action remapping result does not reach significance, which is consistent with
the high variance observed across trials in that condition---some trials received
a near-identity action shuffle while others experienced full remapping---rather
than a conclusive finding of equivalence with the default performance.